\documentclass[11pt]{article}
\usepackage[T1]{fontenc}
\usepackage[utf8]{inputenc}
\usepackage[a4paper,margin=27mm]{geometry}
\usepackage{amsmath,amssymb,amsthm,mathtools,booktabs,microtype}
\usepackage[hidelinks]{hyperref}
\newtheorem{theorem}{Theorem}
\newtheorem{lemma}[theorem]{Lemma}
\newtheorem{proposition}[theorem]{Proposition}
\newtheorem{corollary}[theorem]{Corollary}
\theoremstyle{remark}\newtheorem{remark}[theorem]{Remark}
\newcommand{\CT}{\operatorname{CT}}
\newcommand{\Z}{\mathbb Z}\newcommand{\Q}{\mathbb Q}
\newcommand{\Ztwo}{\mathbb Z_2}\newcommand{\Qtwo}{\mathbb Q_2}
\newcommand{\cO}{\mathcal O}\newcommand{\vp}{v_2}
\newcommand{\vpi}{v_\pi}
\title{Delannoy valuation laws:\\explicit thresholds and multiplicative differences}
\author{C\'esar Augusto Grisa\\{\small Independent Researcher}}
\date{9 November 2026}
\begin{document}
\maketitle
\begin{center}\small
AI-assisted working draft. The mathematical arguments and suggested extensions
require the author's verification; see the assistance statement at the end.
\end{center}
\begin{abstract}
For the central Delannoy numbers $D(N)$, Lengyel conjectured the exact
$2$-adic valuation of $D(u2^{n+1}+b)-D(u2^n+b)$ for odd positive $u$
and sufficiently large $n$. This working revision presents an integrated
argument with explicit thresholds: the valuation is
$n+2v_2(b(b+1))$ for $b\ne0,-1$ and
$n\ge\max\{2,v_2(b(b+1))\}$, provided the indices are nonnegative;
for $b=0,-1$ the valuation is $3n$ for $n\ge2$.
A ramified constant-term argument gives a uniform Mahler bound and
local Taylor estimates. An elementary kernel cancellation gives a cubic
increment congruence from scale $2^3$. The remaining small-valuation
cases reduce to eight explicit arithmetic evaluations involving only
$D(N)$ with $N\le36$. The same ingredients prove the basepoint
valuations in Lengyel's first three conjectures. As a consequence,
for every positive even $n$ and odd integer $p>1$,
$v_2(D(pn)-D(n))=v_2(p-1)+3v_2(n)$.
The finite NEG archive is preserved as supplementary data, while its
claimed total valuation law is derived as a corollary of the analytic proof.
The uniform internal structure of the NEG decomposition is not assumed.
\end{abstract}

\section{The valuation law and its scope}
The central Delannoy numbers, sequence A001850, are given by
\[
D(N)=\sum_{k=0}^N2^k\binom Nk^2\qquad(N\ge0).
\]
Their combinatorial and arithmetic properties are discussed in
\cite{Sulanke2003,BanderierSchwer2005,Sun2018,Lengyel2021}.
Lengyel's Conjecture~4 \cite[Conjecture~4]{Lengyel2021} predicts, for
odd positive $u$ and $n$ sufficiently large, the valuation $3n$ if
$b=0,-1$ and $n+2v_2(b(b+1))$ otherwise.

\begin{theorem}\label{thm:main}
Let $b\in\Z$, let $u$ be an odd positive integer, and let $n$ be a
nonnegative integer. Suppose
$b+u2^n\ge0$.
If $b\notin\{0,-1\}$, put $t=v_2(b(b+1))$. Then
\[
v_2\bigl(D(b+u2^{n+1})-D(b+u2^n)\bigr)=n+2t
\qquad(n\ge\max\{2,t\}).
\]
If $b\in\{0,-1\}$, then
\[
v_2\bigl(D(b+u2^{n+1})-D(b+u2^n)\bigr)=3n
\qquad(n\ge2).
\]
\end{theorem}
The condition on the indices is automatic when $b\ge0$ and eventually
holds for every fixed negative $b$ and positive $u$. Thus the statement
implies the full sufficiently-large-$n$ assertion of Conjecture~4.
It does not assert optimal thresholds in every individual case.
For instance, the generic formula fails at $b=u=1,n=1$:
$D(5)-D(3)=1620$ has valuation $2$, rather than $3$.

\paragraph{Relation to recent work.}
Wang and Zhong \cite[Theorem~3.1]{WangZhong2025} prove
$D(2^a+1)\equiv3+2^{a+2}\pmod{2^{2a+2}}$ for $a\ge2$.
Subtracting consecutive instances already gives the case $b=u=1$
of the generic formula. The target here is the uniform law for arbitrary
integer shifts and odd multipliers, with an explicit threshold.
Recent work on valuations of Legendre polynomial values
\cite{AlekseyevEtAl2025} addresses a different valuation problem;
in particular, its prime-divisibility hypotheses do not directly cover
the present $2$-adic differences at $P_N(3)=D(N)$.

The proof has three parts: analytic control on the four residue classes
modulo $4$, a cubic increment at larger scales, and a finite table for
the remaining two valuation types. The basepoint laws and the fourth conjecture are derived first;
two-point and multiplicative valuation laws follow as corollaries.
The finite NEG archive records a complementary cancellation mechanism.
Historical audit reports and the NEG tables are not mathematical premises
of the universal arguments.

\section{Interpolation and reflection}
Define
\[
F(x)=\sum_{k\ge0}2^k\binom{x}{k}^2\qquad(x\in\Ztwo).
\]
Since $\binom{x}{k}\in\Ztwo$, the series converges uniformly,
defines a continuous function, and agrees with $D$ at nonnegative integers.
\begin{lemma}\label{lem:reflection}
For every $x\in\Ztwo$, one has $F(-x-1)=-F(x)$.
\end{lemma}
\begin{proof}
Euler's hypergeometric transformation gives, for every integer $N\ge0$,
the formal identity
\[
\sum_{k\ge0}\binom{N+k}{k}^2z^k
=(1-z)^{-2N-1}\sum_{j=0}^N\binom Nj^2z^j.
\]
Both sides converge at $z=2$ in $\Qtwo$; expanding the rational factor
about zero is legitimate for $z\in2\Ztwo$. Since
$\binom{-N-1}{k}=(-1)^k\binom{N+k}{k}$, evaluation gives
$F(-N-1)=-D(N)$. Density of the nonnegative integers in $\Ztwo$
and continuity extend the identity to every $x$.
\end{proof}

\section{A uniform Mahler bound on the even charts}
Let $A(u)=F(4u)$ and $H(u)=F(4u+2)$ for $u\in\Ztwo$.
For $u\ge0$ these equal $D(4u)$ and $D(4u+2)$.
Write $\beta_{A,m}=\Delta^m A(0)$ and
$\beta_{H,m}=\Delta^m H(0)$, where $\Delta f(u)=f(u+1)-f(u)$.
We use the usual Mahler expansion of a continuous function on $\Ztwo$
\cite{Mahler1958,Robert2000}. Write $s_2(m)$ for the number of ones
in the binary expansion of $m$.
\subsection{Constant-term model}

Put
\[
\mathcal Q(w)=w+3+2w^{-1}=(1+w)(1+2w^{-1}),
\qquad R(w)=\mathcal Q(w)^4-1.
\]
The elementary constant-term identity
\[
D(n)=\CT\,\mathcal Q(w)^n
\]
follows by expanding the two factors in $(1+w)^n(1+2w^{-1})^n$ and
selecting equal exponents.  Therefore, for every $m\ge1$,
\begin{equation}\label{eq:beta-ct}
\beta_{A,m}=\CT\,R(w)^m,
\qquad
\beta_{H,m}=\CT\,\mathcal Q(w)^2R(w)^m.
\end{equation}
Indeed, applying $\Delta^m$ to $\mathcal Q(w)^{4u}$ at $u=0$ multiplies by
$(\mathcal Q(w)^4-1)^m$.

\subsection{Ramified normalization}

Let $K=\Q_2(\pi)$ with $\pi^2=2$, and let $\cO_K$ be its valuation ring.
Substitute $w=\pi x$ and define
\[
S(x)=\frac{R(\pi x)}4,
\qquad
P(x)=\mathcal Q(\pi x)^2.
\]
A direct expansion gives
\begin{align*}
S(x)={}&x^4+6\pi x^3+31x^2+45\pi x+80\\
&+45\pi x^{-1}+31x^{-2}+6\pi x^{-3}+x^{-4},\\
P(x)={}&13+\pi^2Y+3\pi^3Z,
\end{align*}
where
\[
Y=x^2+x^{-2},\qquad Z=x+x^{-1}.
\]
In particular $S,P\in\cO_K[x,x^{-1}]$ and
\[
S\equiv T\pmod\pi,
\qquad
T=x^4+x^2+x^{-2}+x^{-4}.
\]
Since the substitution $w=\pi x$ does not change a constant term,
\begin{equation}\label{eq:normalization}
\frac{\beta_{A,m}}{4^m}=\CT\,S(x)^m\in\Q_2,
\qquad
\frac{\beta_{H,m}}{4^m}=\CT\,P(x)S(x)^m\in\Q_2.
\end{equation}

\begin{lemma}[Bulk divisibility]\label{lem:bulk}
For every $m\ge1$,
\[
\vp(\beta_{A,m})\ge2m,
\qquad
\vp(\beta_{H,m})\ge2m.
\]
\end{lemma}

\begin{proof}
The two normalized constant terms in \eqref{eq:normalization} belong to
$\cO_K$ because $S$ and $P$ are integral Laurent polynomials.  They also
belong to $\Q_2$.  Hence they are in $\cO_K\cap\Q_2=\Z_2$.  Multiplication
by $4^m$ proves the claim.
\end{proof}

\subsection{One extra factor when $s_2(m)=2$}

\begin{lemma}[Parity of the principal term]\label{lem:parity}
For every $m\ge1$,
\[
\CT\,T^m\equiv0\pmod2.
\]
\end{lemma}

\begin{proof}
Modulo $2$ one has $T\equiv Y^2+Y$.  For every $r\ge1$,
\[
\CT\,Y^r=
\begin{cases}
0,&r\text{ odd},\\
\binom{2a}{a},&r=2a,
\end{cases}
\]
and $\binom{2a}{a}$ is even for $a\ge1$.  Expanding
$(Y^2+Y)^m=\sum_{j=0}^{m}\binom mjY^{m+j}$ proves the result.
\end{proof}

\begin{lemma}[The $s_2(m)=2$ branch]\label{lem:pop2}
If $s_2(m)=2$, then
\[
\vp(\beta_{A,m})\ge2m+1,
\qquad
\vp(\beta_{H,m})\ge2m+1.
\]
\end{lemma}

\begin{proof}
Reduce the normalized terms in \eqref{eq:normalization} modulo $\pi$.
Since $S\equiv T\pmod\pi$ and $P\equiv1\pmod\pi$,
\[
\CT S^m\equiv\CT T^m\equiv0\pmod\pi,
\qquad
\CT(PS^m)\equiv\CT T^m\equiv0\pmod\pi.
\]
Both constant terms are in $\Q_2$.  The intersection
$\pi\cO_K\cap\Q_2=2\Z_2$ now gives one extra factor $2$ beyond
Lemma~\ref{lem:bulk}.
\end{proof}

\subsection{Two extra factors for powers of two}

\begin{lemma}[Power comparison]\label{lem:power-comparison}
For every $a\ge2$,
\[
S^{2^a}\equiv T^{2^a}\pmod{\pi^4\cO_K[x,x^{-1}]}.
\]
\end{lemma}

\begin{proof}
Write $S=T+E$ with $E\in\pi\cO_K[x,x^{-1}]$.  For completeness, the required binomial valuation has a direct proof.
For $1\le k\le2^a-1$,
\[
\binom{2^a}{k}=\frac{2^a}{k}\binom{2^a-1}{k-1}.
\]
Over $\mathbb F_2$,
\[
(1+X)^{2^a-1}=\frac{1+X^{2^a}}{1+X}
=1+X+\cdots+X^{2^a-1},
\]
so $\binom{2^a-1}{k-1}$ is odd.  Hence
\[
\vp\!\binom{2^a}{k}=a-\vp(k).
\]
Thus
\[
\vpi\!\left(\binom{2^a}{k}E^k\right)
\ge2(a-\vp(k))+k\ge4.
\]
For $k=2^a$ the valuation is at least $2^a\ge4$.  The binomial
expansion proves the congruence.
\end{proof}

\begin{lemma}[Parity strengthened at powers of two]\label{lem:power-parity}
For every $a\ge1$,
\[
\CT\,T^{2^a}\in4\Z_2,
\qquad
\CT\,(PT^{2^a})\in4\Z_2.
\]
\end{lemma}

\begin{proof}
Write
\[
T=x^{-4}(1+x^2)(1+x^6).
\]
For $m=2^a$, the constant term is
\[
\CT\,T^m=\sum_b\binom mb\binom m{2m-3b}.
\]
In every nonzero summand both indices are strictly between $0$ and $m$:
the endpoints are invalid, while an index equal to $0$ or $m$ would force
$3\mid m$, impossible for a power of $2$.  Both binomial coefficients
are therefore even, so every summand is divisible by $4$.

Moreover, $\CT(ZT^m)=0$ because $T^m$ has only even exponents, and
\[
\CT(YT^m)=2[x^2]T^m
\]
is even by inversion symmetry.  From
$P=13+\pi^2Y+3\pi^3Z$ and $\pi^2=2$, it follows that
$\CT(PT^m)$ is also divisible by $4$.
\end{proof}

\begin{lemma}[The power-of-two branch]\label{lem:powers}
For every $m=2^a$ with $a\ge0$,
\[
\vp(\beta_{A,m})\ge2m+2,
\qquad
\vp(\beta_{H,m})\ge2m+2.
\]
\end{lemma}

\begin{proof}
For $m=1,2$, direct expansion gives
\[
\CT S=80,\quad \CT(PS)=2244,
\quad \CT S^2=16568,\quad \CT(PS^2)=504968,
\]
all divisible by $4$.  For $m=2^a$ with $a\ge2$,
Lemma~\ref{lem:power-comparison} and Lemma~\ref{lem:power-parity} show
\[
\CT S^m\in4\Z_2,
\qquad
\CT(PS^m)\in4\Z_2.
\]
Multiplying by $4^m$ in \eqref{eq:normalization} proves the claim.
\end{proof}

\subsection{Universal Mahler bound}

\begin{theorem}[Universal Mahler bound]\label{thm:mahler}
For both charts $f=A,H$ and every $m\ge1$,
\[
\vp(\Delta^m f(0))\ge \vp(m!)+m+3.
\]
\end{theorem}

\begin{proof}
Legendre's identity gives $\vp(m!)=m-s_2(m)$, so the desired lower
bound is $2m-s_2(m)+3$.

If $s_2(m)\ge3$, Lemma~\ref{lem:bulk} gives $2m\ge2m-s_2(m)+3$.
If $s_2(m)=2$, Lemma~\ref{lem:pop2} gives $2m+1$.
If $s_2(m)=1$, then $m$ is a power of two and Lemma~\ref{lem:powers}
gives $2m+2$.  These are exactly the three required cases.
\end{proof}

\subsection{Mahler-to-Tate-to-Taylor consequence}

\begin{corollary}[Tate and Taylor control]\label{cor:tate}
For $f=A,H$, there is a series $G_f(u)\in\Z_2\langle u\rangle$ such that
\[
f(u)-f(0)=16uG_f(u).
\]
Moreover,
\[
f(u+h)-f(u)-hf'(u)\in32h^2\Z_2\langle u,h\rangle.
\]
\end{corollary}

\begin{proof}
Write $\beta_m=\Delta^m f(0)$ and
\[
f(u)=\sum_{m\ge0}\beta_m\binom um
=\sum_{m\ge0}c_m(u)_m,
\qquad c_m=\frac{\beta_m}{m!}.
\]
Theorem~\ref{thm:mahler} gives $c_m\in2^{m+3}\Z_2$ for $m\ge1$.
If
\[
(u)_m=\sum_{j=0}^m s(m,j)u^j
\]
with $s(m,j)\in\Z$, then the power-series coefficient
\[
a_j=\sum_{m\ge j}c_ms(m,j)
\]
converges in $\Z_2$ and satisfies $a_j\in2^{j+3}\Z_2$ for $j\ge1$.
Hence $f(u)-f(0)=\sum_{j\ge1}a_ju^j$ belongs to
$16u\Z_2\langle u\rangle$.

Finally,
\[
f(u+h)-f(u)-hf'(u)
=\sum_{j\ge2}a_j\sum_{\ell=2}^j\binom j\ell u^{j-\ell}h^\ell.
\]
For $j\ge2$, $a_j\in2^{j+3}\Z_2\subseteq32\Z_2$.  Factoring $h^2$
proves the asserted quadratic remainder bound.
\end{proof}


\section{Local estimates and eight arithmetic evaluations}
% Self-contained replacement for M3 -> M5a -> M5b assembly.
% Inputs: reflection and a_j in 2^{j+3} Z_2 for A and H, proved in M3.
% Mathematical improvement independently derived 20 October 2026.

\begin{lemma}[Local analytic control]\label{lem:local-control}
For each $r\in\{0,1,2,3\}$ write
\[
G_r(u)=F(4u+r)=\sum_{j\ge0}a_{r,j}u^j.
\]
Then $a_{r,j}\in2^{j+3}\mathbb Z_2$ for $j\ge1$. Consequently, for
$b\in\mathbb Z_2$ and $h,L\in4\mathbb Z_2$,
\begin{align}
F(b+h)-F(b)-hF'(b)&\in2h^2\mathbb Z_2,\label{eq:physical-taylor}\\
F'(b+h)-F'(b)&\in4h\mathbb Z_2,\label{eq:physical-derivative}\\
F(b+h+L)-F(b+h)-F(b+L)+F(b)&\in4hL\mathbb Z_2.\label{eq:physical-mixed}
\end{align}
In particular, if $\Delta(b,h)=F(b+2h)-F(b+h)$, then
\[
\Delta(b,h)-hF'(b)\in2h^2\mathbb Z_2.
\]
\end{lemma}
\begin{proof}
The coefficient bound for $r=0,2$ is the Mahler--Tate conclusion. For the
other two charts use $G_1(u)=-G_2(-u-1)$ and
$G_3(u)=-G_0(-u-1)$. The coefficient of $u^j$ after this substitution is
a convergent sum of terms $\pm a_{r,k}\binom{k}{j}$, $k\ge j$;
every such term lies in $2^{j+3}\mathbb Z_2$.

For any of the four charts, the ordinary Taylor remainder lies in
$32v^2\mathbb Z_2\langle u,v\rangle$, since its terms have $j\ge2$.
Also
\[
G_r'(u+v)-G_r'(u)\in64v\mathbb Z_2\langle u,v\rangle:
\]
the $j=1$ contribution vanishes, the $j=2$ contribution has the extra
factor $2$, and $a_{r,j}\in64\mathbb Z_2$ for $j\ge3$.
The same argument gives
\[
G_r(u+v+w)-G_r(u+v)-G_r(u+w)+G_r(u)
\in64vw\mathbb Z_2\langle u,v,w\rangle.
\]
Indeed, its degree-two contribution is $2a_{r,2}vw$; all higher-degree
contributions have integral coefficients and $a_{r,j}\in64\mathbb Z_2$.
Writing $b=4u+r$, $h=4v$, and $L=4w$ gives
\eqref{eq:physical-taylor}--\eqref{eq:physical-mixed} after division by
$4$ for the derivative or by $16$ for the quadratic factors.
Finally, apply \eqref{eq:physical-taylor} to the increments $2h$ and $h$
at $b$ and subtract.
\end{proof}

\begin{lemma}[Four derivative checks]\label{lem:four-derivatives}
Let $b\in\Ztwo$. If $t=v_2(b(b+1))\in\{1,2\}$, then
\[
F'(b)\equiv2^{2t}\pmod{2^{2t+1}}.
\]
\end{lemma}
\begin{proof}
Put $H=2^{2t}$ and $M=2^{2t+1}$. Lemma~\ref{lem:local-control} gives
\[
\frac{F(b+2H)-F(b+H)}{H}\equiv F'(b)\pmod{M}.
\]
For $t=1$, the relevant classes are $b=1,2\pmod4$; derivative shifts
by $4$ lie in $16\mathbb Z_2\subseteq M\mathbb Z_2$.
For $t=2$, the relevant classes are $b=3,4\pmod8$; derivative shifts
by $8$ lie in $32\mathbb Z_2=M\mathbb Z_2$.
It therefore suffices to evaluate the following four entries from
$D(N)=\sum_{k=0}^{N}2^k\binom Nk^2$:
\[
\begin{array}{c|c|c|c|c}
t&b&H&M&[D(b+2H)-D(b+H)]/H\pmod M\\\hline
1&1&4&8&4\\
1&2&4&8&4\\
2&3&16&32&16\\
2&4&16&32&16
\end{array}
\]
All four quotients are integers and the listed residues are exact.
\end{proof}

\begin{lemma}[Four transition checks]\label{lem:four-transitions}
Let $b\in\Ztwo$. If $t=v_2(b(b+1))=2$ and $n\in\{2,3\}$, then for every odd $u$,
\[
v_2\bigl(\Delta(b,u2^n)\bigr)=n+4.
\]
\end{lemma}
\begin{proof}
Work modulo $2^{n+5}$. By \eqref{eq:physical-mixed}, replacing $b$ by
$b+8$ changes $\Delta(b,h)$ by a multiple of $4h\cdot8$, which is
divisible by $2^{n+5}$ when $v_2(h)=n$.

To reduce $u$ modulo $2$, set $K=2^{n+1}$. Taylor expansion at $b+2h$
and $b+h$ gives
\[
\Delta(b,h+K)-\Delta(b,h)
=K\bigl(2F'(b+2h)-F'(b+h)\bigr)+E,
\quad E\in2K^2\mathbb Z_2.
\]
By \eqref{eq:physical-derivative}, the expression in parentheses belongs
to $F'(b)+4h\mathbb Z_2\subseteq16\mathbb Z_2$, using
Lemma~\ref{lem:four-derivatives} and $n\ge2$. Its product with $K$ is
therefore divisible by $2^{n+5}$; so is $E$, since
$v_2(2K^2)=2n+3\ge n+5$.
Thus $b$ reduces to $3,4\pmod8$ and odd $u$ reduces to $1\pmod2$.
The four exact evaluations are
\[
\begin{array}{c|c|c|c|c}
n&b&u&\text{modulus}&D(b+2u2^n)-D(b+u2^n)\\\hline
2&3&1&128&64\pmod{128}\\
2&4&1&128&64\pmod{128}\\
3&3&1&256&128\pmod{256}\\
3&4&1&256&128\pmod{256}
\end{array}
\]
These residues have valuation $n+4$ as required.
\end{proof}

\begin{proposition}[Low-valuation branch]\label{prop:low-eight-checks}
For $b\in\Ztwo$ with $t=v_2(b(b+1))\in\{1,2\}$ and every odd $u$, the equality
\[
v_2\bigl(F(b+2u2^n)-F(b+u2^n)\bigr)=n+2t
\]
holds for every $n\ge2$.
\end{proposition}
\begin{proof}
For $n\ge2t$, set $h=u2^n$. The main term $hF'(b)$ has valuation
$n+2t$, by Lemma~\ref{lem:four-derivatives}. Its error lies in
$2h^2\mathbb Z_2$ and has valuation at least $2n+1\ge n+2t+1$.
The sum therefore has valuation $n+2t$. This covers every $n\ge2$ for
$t=1$ and every $n\ge4$ for $t=2$. The two remaining values of $n$ are
Lemma~\ref{lem:four-transitions}.
\end{proof}

\section{An elementary cubic increment}
% Self-contained replacement for the original M2/M4 arguments.
% Uses only amsmath, amssymb, amsthm and ordinary theorem environments.

\begin{lemma}[Elementary kernel cancellation]\label{lem:elementary-kernels}
Put
\[
 H_j=\sum_{i=1}^j\frac1i,\qquad H_0=0,\qquad
 K=\sum_{k\ge1}\frac{2^k}{k^2},\qquad
 \Sigma=\sum_{k\ge1}\frac{2^kH_{k-1}}{k^2},\qquad\nu=-2\Sigma.
\]
These series converge in $\mathbb Q_2$, and
$K=0$ and $\nu\in\mathbb Z_2^\times$.
\end{lemma}
\begin{proof}
Write
$L(z)=\sum_{k\ge1}z^k/k^2$ and
$\ell(z)=\sum_{k\ge1}(-1)^{k+1}z^k/k$.
Both series converge when $z\in2\mathbb Z_2$.
The following identity holds in $\mathbb Q[[z]]$:
\[
 L(z)+L\left(\frac{z}{z-1}\right)
 =-\frac12\ell(-z)^2.
\]
Indeed, differentiation makes both sides equal to
$\ell(-z)/(1-z)$, and both sides have zero constant term.
All series and the composition converge for $z\in2\mathbb Z_2$.
For example, expand $(z/(z-1))^k=(-1)^kz^k(1-z)^{-k}$;
its degree-$d$ terms after division by $k^2$ have valuation at least
$d-2\lfloor\log_2d\rfloor$ for $d\ge k$.
This tends to infinity, justifying the double-series rearrangement.
Thus the formal identity can be evaluated at $z=2$.
The formal logarithm identity
\[
 \ell(a+b+ab)=\ell(a)+\ell(b)
\]
also converges for $a,b\in2\mathbb Z_2$.  At $a=b=-2$ it gives
$0=2\ell(-2)$, hence $\ell(-2)=0$ in characteristic zero.
Since $2/(2-1)=2$, the first identity now gives $2K=0$.

For the second assertion, set
$T_k=2^{k+1}H_{k-1}/k^2$.
For $k\ge2$,
\[
 v_2(T_k)\ge k+1-2v_2(k)-\lfloor\log_2(k-1)\rfloor.
\]
This bound tends to infinity, proving convergence.
For $k=2,3,5,6,7,8$ the bound is at least one, whereas
$T_4=11/3$ is odd.  For $k\ge9$, put
$a=\lfloor\log_2 k\rfloor$; the bound is at least $k+1-3a$,
which is at least one for $a=3$ and at least
$2^a+1-3a\ge5$ for $a\ge4$.
Thus $\nu=-\sum_kT_k\equiv1\pmod2$.
\end{proof}

\begin{proposition}[Cubic increment from scale eight]\label{prop:cubic-t3}
Let
$F(z)=\sum_{k\ge0}2^k\binom zk^2$ for $z\in\mathbb Z_2$.
For integers $x\ne y$ and $t\ge3$, put
$q=2^t$ and $c=v_2(y-x)$.  Then
\[
 F(qy)-F(qx)
 \equiv\nu q^3(y^3-x^3)\pmod{2^{3t+c+1}}.
\]
\end{proposition}
\begin{proof}
Use the cutoff $Q=2^{t+1}$.
For $2\le k<Q$, write
\[
 R_k(X)=\prod_{j=1}^{k-1}(1-X/j)^2
        =\sum_{m=0}^{2k-2}e_{k,m}X^m,
 \qquad L_k=\lfloor\log_2(k-1)\rfloor.
\]
Each factor has degree-$d$ coefficient of valuation at least $-dL_k$;
hence $v_2(e_{k,m})\ge-mL_k$.
Also $e_{k,0}=1$, $e_{k,1}=-2H_{k-1}$ and
\[
 2^k\binom{qz}{k}^2=q^2z^2\frac{2^k}{k^2}R_k(qz).
\]
For every $t\ge3$ and $k\ge2$,
\[
 t+k-2v_2(k)-2L_k\ge1.                         \tag{*}
\]
For $L_k=0,1,2$ the minimum is $t-2$.
If $L=L_k\ge3$ and $2^L<k<2^{L+1}$, then
$v_2(k)\le L-1$, giving the lower bound
$t+2^L+3-4L\ge t-1\ge2$.
If $k=2^{L+1}$, the expression is
$t+2^{L+1}-4L-2\ge t+2$.

Since $L_k\le t$ for $k<Q$, the degree-$m$ part with $m\ge2$
of the difference has valuation at least
\[
 3t+c+[t+k-2v_2(k)-2L_k]+(m-2)(t-L_k)
 \ge3t+c+1.
\]
Here $y^{m+2}-x^{m+2}$ is divisible by $y-x$.
The $k=1$ term is exactly $2q^2(y^2-x^2)$, so including it yields
\[
 \sum_{1\le k<Q}2^k
  \left[\binom{qy}{k}^2-\binom{qx}{k}^2\right]
 \equiv q^2K_{<Q}(y^2-x^2)
       -2q^3\Sigma_{<Q}(y^3-x^3)
       \pmod{2^{3t+c+1}},
\]
where $K_{<Q}=\sum_{1\le k<Q}2^k/k^2$ and
$\Sigma_{<Q}=\sum_{1\le k<Q}2^kH_{k-1}/k^2$.

For the quadratic kernel tail, with $a=\lfloor\log_2k\rfloor\ge t+1$,
\[
 v_2(q^2 2^k/k^2)\ge2t+2^a-2a\ge Q-2.
\]
The last expression increases with $a$ since its successive difference
is $2^a-2>0$.
For the cubic kernel tail, at $k=2^a$ the elementary bound is
$3t+2+2^a-3a\ge Q-1$; between consecutive powers of two
it is at least two larger.  Again the dyadic minimum increases,
as its successive difference is $2^a-3>0$.
Consequently
\[
 q^2(K-K_{<Q})\in2^{Q-2}\mathbb Z_2,
 \qquad
 2q^3(\Sigma-\Sigma_{<Q})\in2^{Q-1}\mathbb Z_2.
\]
After multiplication by the respective power differences, each bound
gains $c$.  Since $Q-2=2^{t+1}-2\ge3t+1$ for $t\ge3$,
the finite kernels can be replaced by $K$ and $\Sigma$.

Finally, generalized binomial coefficients at integer arguments are
integers, including negative arguments.  Vandermonde's identity and
$\binom{q(y-x)}j=\frac{q(y-x)}j\binom{q(y-x)-1}{j-1}$ give
\[
 v_2\left(2^k\left[\binom{qy}{k}^2-\binom{qx}{k}^2\right]\right)
 \ge k+t+c-\lfloor\log_2k\rfloor
 \ge Q+c-1\qquad(k\ge Q).
\]
The function $k-\lfloor\log_2k\rfloor$ is nondecreasing.
The high-index sum therefore lies in
$2^{Q+c-1}\mathbb Z_2\subseteq2^{3t+c+1}\mathbb Z_2$.
Using $K=0$ and $\nu=-2\Sigma$ finishes the proof.
\end{proof}

\begin{corollary}[Odd-base valuation]\label{cor:odd-base-t3}
For $t\ge3$, odd $r\in\mathbb Z$, and nonzero $s\in\mathbb Z$,
\[
 v_2\left(F(2^t(r+2s))-F(2^t(r+s))\right)=3t+v_2(s).
\]
\end{corollary}
\begin{proof}
The cubic difference is
$s(3r^2+9rs+7s^2)$, whose parenthesis is odd for either parity of $s$.
Since $\nu$ is a unit, the main term has valuation $3t+v_2(s)$
and the error has strictly larger valuation.
\end{proof}


\section{Basepoint valuations and the first three conjectures}
\begin{proposition}[Valuations at zero and minus one]\label{prop:basepoints}
For every odd integer $u$ and every $a\ge2$,
\[
v_2\bigl(F(2^a u)-1\bigr)=3a,
\qquad
v_2\bigl(F(2^a u-1)+1\bigr)=3a.
\]
For positive $u$, the same statements hold with $D$ in place of $F$.
\end{proposition}
\begin{proof}
For $a\ge3$, apply Proposition~\ref{prop:cubic-t3} with $t=a$,
$x=0$, and $y=u$. Since $F(0)=1$, $v_2(u)=0$, and $\nu$ is a unit,
\[
F(2^a u)-1\equiv\nu2^{3a}u^3\pmod{2^{3a+1}},
\]
which has valuation $3a$. This allows either sign of $u$.

It remains to treat $a=2$. For every odd integer $u$, put $b=4u$.
Then $v_2(b(b+1))=2$, and Lemma~\ref{lem:four-derivatives} gives
$F'(b)\in16\Ztwo$. By \eqref{eq:physical-taylor} with $h=8$,
\[
F(b+8)-F(b)=8F'(b)+E,\qquad E\in128\Ztwo.
\]
Thus $F(b+8)\equiv F(b)\pmod{128}$. Every odd integer $u$ has
the form $1+2j$ with $j\in\Z$; iterating this congruence in either
direction gives
\[
F(4u)\equiv F(4)=321\pmod{128}.
\]
Consequently $F(4u)-1\equiv64\pmod{128}$, of valuation $6$.

For the second equality, reflection yields
\[
F(2^a u-1)+1=-\bigl(F(-2^a u)-1\bigr).
\]
The first part applies to the signed odd integer $-u$.
At nonnegative indices $F$ agrees with $D$.
\end{proof}

\begin{corollary}[Lengyel's first three conjectures]\label{cor:c123}
For $a\ge2$ and positive odd $u$,
\[
v_2\bigl(D(2^a u)-1\bigr)
=v_2\bigl(D(2^a u-1)+1\bigr)=3a.
\]
In particular,
\[
v_2\bigl(D(4m+3)+1\bigr)=3\bigl(v_2(m+1)+2\bigr)
\qquad(m\ge0),
\]
and
\[
v_2\bigl(D(2m)-1\bigr)=3v_2(2m)
\qquad(m>0\text{ even}).
\]
\end{corollary}
\begin{proof}
The first statement is Proposition~\ref{prop:basepoints} and is
Conjecture~1 of \cite{Lengyel2021}. For Conjecture~2, apply the
minus-one formula to $4(m+1)$, whose valuation is $2+v_2(m+1)$.
For Conjecture~3, use the zero-basepoint formula at $2m$.
The degenerate value $m=0$ in the last formula can be included only
with the extended convention $v_2(0)=+\infty$.
\end{proof}

\section{Assembly and the exceptional shifts}
\begin{proof}[Proof of Theorem~\ref{thm:main}]
First suppose $b\notin\{0,-1\}$, and put $t=v_2(b(b+1))$.
For $t=1,2$, Proposition~\ref{prop:low-eight-checks} proves the claim
for $F$ on $\Ztwo$. For $t\ge3$, set $q=2^t$ and
$s=u2^{n-t}$, so $s$ is a nonzero integer.
If $b=qr$, then $r$ is odd, and Corollary~\ref{cor:odd-base-t3} gives
\[
v_2\bigl(F(b+2u2^n)-F(b+u2^n)\bigr)
=3t+v_2(s)=n+2t.
\]
If $b=qr-1$ with $r$ odd, reflection turns the difference into
\[
-\bigl(F(q(-r-2s))-F(q(-r-s))\bigr).
\]
Apply the same corollary with $(-r,-s)$. The argument allows negative
$r$ and therefore covers negative integer shifts as well.

For $b=0$ and $n\ge2$, Proposition~\ref{prop:basepoints} gives
\[
v_2\bigl(D(u2^{n+1})-1\bigr)=3n+3,
\qquad
v_2\bigl(D(u2^n)-1\bigr)=3n.
\]
Subtracting these deviations from $1$ gives valuation $3n$.
For $b=-1$, use the two deviations from $-1$ furnished by the same
proposition; their valuations are again $3n+3$ and $3n$.
In each case, adding an error of strictly higher valuation preserves
the valuation of the main term: $B+E=B(1+E/B)$ with $E/B\in2\Ztwo$.
Finally $F$ agrees with $D$ at both indices, which are nonnegative
by the hypothesis of the theorem.
\end{proof}


\section{Two-point and multiplicative valuation laws}
\begin{corollary}[Two even indices of the same valuation]\label{cor:two-point}
Let $M,N$ be distinct positive even integers with
$v_2(M)=v_2(N)=a$. Then
\[
v_2\bigl(D(M)-D(N)\bigr)=v_2(M-N)+2a.
\]
\end{corollary}
\begin{proof}
Exchange $M,N$ if necessary so that $M>N$. Put
$h=M-N=U2^k$ with $U$ positive odd and $k=v_2(M-N)$.
Since $M/2^a$ and $N/2^a$ are odd, $k\ge a+1\ge2$.
Set $B=2N-M$. Then $v_2(B)=a$, $B+1$ is odd, and
$B\notin\{0,-1\}$. In particular $v_2(B(B+1))=a$.
Moreover
\[
B+U2^k=N>0,\qquad B+U2^{k+1}=M>0.
\]
Theorem~\ref{thm:main} applies because $k\ge\max\{2,a\}$ and gives
the required valuation $k+2a$.
\end{proof}

\begin{corollary}[Global even multiplier law]\label{cor:multiplier}
For every positive even integer $n$ and every odd integer $p>1$,
\[
\boxed{v_2\bigl(D(pn)-D(n)\bigr)=v_2(p-1)+3v_2(n).}
\]
No primality assumption on $p$ is required.
\end{corollary}
\begin{proof}
Write $a=v_2(n)$ and $s=v_2(p-1)$. The two indices $pn,n$ have
the same valuation $a$, and
$v_2(pn-n)=a+s$. Corollary~\ref{cor:two-point} gives
$(a+s)+2a=s+3a$.
\end{proof}

\begin{remark}[The explicit bridge to the NEG target]\label{rem:neg-bridge}
Writing $n=2^a u$ and $p-1=2^s q$ with $u,q$ odd, the direct
substitution in Theorem~\ref{thm:main} is
\[
B=(2-p)n,\qquad k=a+s,\qquad U=qu.
\]
Here $B+U2^k=n$, $B+U2^{k+1}=pn$, and
$v_2(B(B+1))=a$. Since $a,s\ge1$, $k\ge\max\{2,a\}$.
The negative-shift part of the theorem is essential when $p>1$.
An unspecified sufficiently-large-scale assertion would not justify
this substitution uniformly, because $B$ varies with the parameters.
For example, $p=3,n=2$ gives
$D(6)-D(2)=8976=2^4\cdot561$, consistent with $1+3=4$.
\end{remark}

\section{Finite arithmetic and reproducibility}
The only finite evaluations required by the proof are the four derivative
residues and four transition residues displayed in the two tables.
They are direct evaluations of a finite binomial sum and involve indices
at most $36$. Thus the written proof can be checked without a large
computation or access to an external archive.
The accompanying file \texttt{eight\_checks.py} recomputes all eight
entries using exact integers. A separate modular script tests the
cubic congruence for $8{,}832$ signed pairs and scales $3\le t\le8$;
these additional checks corroborate the estimates and are not substitutes
for their universal proofs. The original million-class transition
certificate is no longer needed by this revision.


\section{The archived finite NEG data}\label{sec:neg-archive}
The accompanying archive \texttt{NEG\_original.zip} preserves the
submitted finite dataset and its reports. Its target formula is exactly
Corollary~\ref{cor:multiplier}. The archived section decomposes a
truncated sum into terms with negative and nonnegative valuation margins
relative to $w=v_2(p-1)+3v_2(n)$. In this terminology, NEG refers to
negative margins of summands, not to negative Delannoy indices.

The report describes $104{,}448$ sampled parameter tuples and
$13{,}552$ NEG cases. Its CSV and JSON contain $11{,}858$ distinct
keys of the form
\[
(p\bmod8,\ s,\ a,\ q\bmod256,\ u\bmod128).
\]
An archive-consistency check confirms agreement between those two tables,
the NEG multiplicity total, the stored ladder lengths, and the terminal
pair $(-1,0)$ in every recorded ladder. The full sampling total is
reported by the archive; only the NEG state records are supplied.
The archived generation parameters have $s\le12$ and a finite list
of exponents $a\le64$. The actual NEG table has $a\ge7$.

The package does not include the program that generated the terms,
their complete formulas, a proof of reduction to all residue states,
or a written proof of its claimed tail bound. In particular, the key
retains the unbounded parameters $a,s$ exactly. Internal consistency
of the stored ladders is therefore distinct from a universal proof of
their dependence on the key.

The archived summands are not identified with the harmonic-kernel terms
defined earlier. Some cumulative valuations rise and then fall: for
example, a stored ladder contains $(-3,1),(-2,-1),(-1,0)$. The terminal
pair therefore does not assert monotonicity.

The role of this archive in the unified paper is supplementary:
Corollary~\ref{cor:multiplier} proves the total valuation law for all
the stated even indices and odd multipliers, including parameters
outside the finite sample, without invoking the NEG decomposition.
It does not by itself prove the separate valuation of each block,
the determinism of every unobserved ladder, or the archived tail claim.
Those internal statements require their own definitions and proofs.

\paragraph{Reproduction files.}
\texttt{check\_joined\_claims.py} tests the added basepoint,
exceptional, two-point, and multiplier consequences by exact integer
arithmetic and tests additional large parameters modularly.
\texttt{verify\_NEG\_archive.py} checks only archive integrity and
cross-table consistency. Its JSON report explicitly distinguishes
these checks from regeneration of the original mathematical computation.

\paragraph{Assistance and draft status.}
This working revision was prepared with OpenAI Codex assistance for
mathematical review, suggested proof techniques, exact arithmetic,
literature search, and English drafting. The proposed changes include
the doubled cutoff in the cubic estimate, the stronger local Taylor
and mixed-increment bounds, the eight-entry reduction, the elementary
kernel proof, the extensions to negative and exceptional shifts,
the basepoint argument at scale four, the deductions of Conjectures~1--3,
and the two-point and multiplier corollaries. The finite NEG archive
was checked for internal consistency, not regenerated from missing source.
The author must independently verify and take responsibility for these
arguments and references. This file is a review draft, not a claim of
completed peer review. Before submission, its narrative must be reconstructed
by the author where required by the chosen journal's policy, and the
assistance must be disclosed as that policy requires.

\begin{thebibliography}{10}
\bibitem{AlekseyevEtAl2025}
M.~A. Alekseyev, T. Amdeberhan, J. Shallit, and I. Vukusic,
\emph{On the $p$-adic valuations of values of Legendre polynomials},
arXiv:2505.08935 (2025),
\url{https://arxiv.org/abs/2505.08935}.
\bibitem{BanderierSchwer2005}
C. Banderier and S. Schwer, Why Delannoy numbers?,
\emph{J. Statist. Plann. Inference} \textbf{135} (2005), 40--54.
\bibitem{Lengyel2021}
T. Lengyel, On some $p$-adic properties and supercongruences of Delannoy
and Schr\"oder numbers, \emph{INTEGERS} \textbf{21} (2021), A86,
\url{https://math.colgate.edu/~integers/v86/v86.pdf}.
\bibitem{Mahler1958}
K. Mahler, An interpolation series for continuous functions of a $p$-adic
variable, \emph{J. Reine Angew. Math.} \textbf{199} (1958), 23--34;
correction, \textbf{208} (1961), 70--72.
\bibitem{Robert2000}
A.~M. Robert, \emph{A Course in $p$-adic Analysis},
Graduate Texts in Mathematics 198, Springer, 2000.
\bibitem{Sulanke2003}
R.~A. Sulanke, Objects counted by the central Delannoy numbers,
\emph{J. Integer Seq.} \textbf{6} (2003), Article 03.1.5.
\bibitem{Sun2018}
Z.-W. Sun, Arithmetic properties of Delannoy numbers and Schr\"oder numbers,
\emph{J. Number Theory} \textbf{183} (2018), 146--171.
\bibitem{WangZhong2025}
R.-H. Wang and M.~X.~X. Zhong,
\emph{Congruences for sums of Delannoy numbers and polynomials},
arXiv:2505.05728v1 (2025), Theorem~3.1,
\url{https://arxiv.org/html/2505.05728v1}.
Published under the title \emph{Congruences for sums of Delannoy numbers
and Legendre polynomials}, \url{https://doi.org/10.1142/S1793042126500442}.
The theorem numbering cited here refers to the consulted preprint.
\end{thebibliography}
\end{document}
